\section{Introduction}

\subsection{The portal}

Bhoonidhi (\url{https://bhoonidhi.nrsc.gov.in/}) is the National Remote
Sensing Centre's public repository for Indian Earth observation data. It
holds satellite imagery, thematic maps and derived products across
approximately 59 published collections spanning the ResourceSat, Cartosat,
Oceansat, EOS and RISAT missions, together with mirrored international
holdings such as Sentinel-1 and Sentinel-2.

Alongside the browser portal, NRSC publishes an \emph{Open Data Access
API}~\cite{bhoonidhiapi}. It is worth stating plainly at the outset what that
API does well, because an argument that misrepresents its target is worthless:

\begin{itemize}
  \item it is built on \textbf{STAC 1.0.0}~\cite{stac}, a current community
        standard for Earth observation metadata;
  \item it has proper token authentication (\ep{/auth/token},
        \ep{/auth/logout});
  \item \ep{/data/collections} enumerates the published collections;
  \item \ep{/data/search} accepts a GeoJSON \code{intersects} geometry, a
        \code{bbox}, RFC 3339 \code{datetime} ranges and CQL2-JSON filters;
  \item \ep{/download} delivers a product by item identifier.
\end{itemize}

That is a competent discovery API. A developer can authenticate, search by
space and time, and fetch. \textbf{Discovery is solved, and this project
claims no credit for noticing otherwise.}

\subsection{The task}

The practical brief asks for one thing: identify a dataset or data product on
Bhoonidhi that would genuinely benefit from an OGC service interface which
does not currently exist, verify that the gap is real, then build and deploy
a conformant service that fills it. Gap identification carries 5 of the 20
marks and is assessed on whether the justification is \emph{convincing}; OGC
compliance carries 10; a working deployment carries 5.

\subsection{What this project identifies, and what it builds}

The gap identified here is not in discovery and not in access control. It is
in \textbf{computation}.

Bhoonidhi stores the footprint polygon of every scene it holds. Its whole
spatial surface, audited against its own published specification, is
\code{bbox} and \code{intersects} --- which are one predicate expressed two
ways: \emph{does this item intersect my area of interest?} There is no union,
no difference, no area, no aggregation, no statistics, and no processing
endpoint anywhere in ISRO's public stack.

The practical consequence is that a user cannot learn \emph{which} of the
scenes a search returns they actually need. Lacking that, the rational
behaviour --- and the observed behaviour --- is to order everything the search
returned. Section~\ref{sec:cost} measures what that costs: across eight
queries the catalogue returns 208 scenes where 16 cover the same ground to
the same percentage.

This report documents a conformant \textbf{OGC API --- Processes}
service~\cite{ogcprocesses} that closes the gap, exposing six operations:

\begin{center}
\small
\begin{tabularx}{\textwidth}{@{}l>{\raggedright\arraybackslash}X@{}}
\toprule
\textbf{Process} & \textbf{The question it answers} \\
\midrule
\code{aoi-coverage}         & How much of my AOI do these scenes cover, and where are the holes? \\
\code{minimal-scene-set}    & Which few of the scenes returned do I actually need? \\
\code{multi-sensor-overlap} & Where do optical and radar co-cover, for fusion? \\
\code{revisit-density}      & How uneven is acquisition density across my AOI? \\
\code{zonal-statistics}     & What are the raster values inside this polygon? \\
\code{scene-raster-profile} & What terrain is actually underneath each candidate scene? \\
\bottomrule
\end{tabularx}
\end{center}

The first four read only footprint geometry --- deciding which scenes are
worth fetching. The last two read pixels --- establishing what is actually
inside them. \code{scene-raster-profile} joins the two tiers. This is
deliberately \emph{not} spectral index computation: there is no band maths and
no derived imagery product, because the gap being filled is spatial algebra
over the catalogue's own geometry.

A supporting \textbf{OGC API --- Records}~\cite{ogcrecords} catalogue of
Bhoonidhi's collections ships alongside the Processes tier. It is
infrastructure rather than the gap being filled: the spatial processes need
discoverable collection identifiers to point at. The substance of the
submission rests on the Processes tier.

\subsection{How to read this report}

Section~\ref{sec:methodology} states how the gap was established and what
counts as evidence. Section~\ref{sec:gap} is the gap analysis and is the core
of the argument. Sections~\ref{sec:design} and~\ref{sec:implementation} cover
service design and implementation; Section~\ref{sec:compliance} the OGC
compliance evidence; Section~\ref{sec:results} the live demonstration;
Section~\ref{sec:limitations} the limitations, stated plainly. Every number in
this report is produced by a script in the repository and saved under
\code{docs/evidence/}; Appendix~\ref{app:reproducing} lists the commands.
